\section{Discussion}
\label{sec:discussion}

\paragraph{Why the open library fails.} A big laboratory library is a strength in the lab and a liability on Mars. Among \nSpecies{} species there are always some, mostly glasses and other amorphous phases, whose broad bands fit SuperCam's broad fiber bump better than any crystalline mineral fits the weak, broadened Mars peaks. Nearest-neighbor matching has no way of knowing that the broadest feature in a Mars spectrum comes from the instrument. Two kinds of knowledge that mission scientists apply without thinking fix most of this: which minerals are plausible at the site, and which parts of the spectrum to trust. Encoding both explicitly took library search from \balCosineOpen{} to \balCosineWindowMars{} balanced accuracy.

\paragraph{The skip has a price.} Skipping the region below 560\,cm$^{-1}$ also throws away real mineral bands, such as quartz near 464\,cm$^{-1}$, feldspar near 480--510\,cm$^{-1}$, and the secondary anhydrite bands near 417 and 499\,cm$^{-1}$, all of which appear in the team's table. This shows up in the per-group recalls: skipping the region raises olivine recall from \recOlivineCosineMars{} to \recOlivineCosineWindowMars{} but lowers sulfate recall from \recSulfateCosineMars{} to \recSulfateCosineWindowMars{}, as anhydrite becomes easier to confuse with pyroxene in the 1000--1020\,cm$^{-1}$ region. A better treatment would model the fiber bump rather than delete the region. Our fiber-bump template is one attempt, and it helped the CNN (\balCnnMarsifiedArtifactMars{} vs.\ \balCnnMarsifiedMars{}) without solving the problem.

\paragraph{Simulated robustness is not real robustness.} The most important negative result of this study is that our simulated-Mars training worked as intended in the lab and hurt on Mars. That training is a reasonable recipe: blurring to the instrument's resolution, calibration shifts, mixing, smooth backgrounds, and noise. Two lessons follow. First, a degradation model is a hypothesis about the instrument, and a wrong hypothesis can teach a model to rely on the wrong cues. For example, a CNN trained on smooth synthetic backgrounds may learn to ignore exactly the kind of broad structure that the fiber bump adds, then misread it at test time. Second, evaluating Mars-oriented classifiers only on synthetically degraded lab spectra, as is often done, can give a misleading picture. Labeled flight spectra are the only way to check. This is the main reason we release the benchmark.

\paragraph{Simple and transparent does well.} The highest-scoring method in this study is also the simplest: cosine matching against a lab library, with a physically motivated restriction and mask. This agrees with earlier lab results in which full-spectrum similarity matching did as well as or better than trained classifiers \citep{carey2015machine}. It is also practical for mission use, because every decision can be traced to a specific reference spectrum. Deep models may still win with better training data, for example spectra from SuperCam-like laboratory systems \citep{manrique2024simulcam}, or with unsupervised domain adaptation on the thousands of unlabeled SuperCam spectra in the archive \citep{ganin2016domain}. Our results suggest that such methods should be judged on flight data, not on simulations.

\subsection{Limitations}
\label{sec:limits}
\paragraph{The labels are expert interpretations.} Our reference labels are the SuperCam team's identifications, not independent ground truth. Some are marked low confidence, and some lean on SuperCam's other techniques (LIBS and visible/infrared spectroscopy) as well as Raman. A model that ``disagrees'' with a low-confidence label is not necessarily wrong, which is why we also report results by confidence level (Section~\ref{sec:results}).

\paragraph{Only spectra with detectable minerals are included.} The team's table lists only spectra that contain identifiable mineral features \citep{jpl2026supercamsis}. Our test therefore asks ``\emph{which} mineral is this?'' and never ``is there a mineral signal here at all?''. A deployed system would also need to reject spectra with no usable signal, and our numbers say nothing about that.

\paragraph{Small, imbalanced, clustered data.} The test set has \nMarsEval{} spectra from only \nMarsTargets{} rocks. Olivine dominates it, and silica, feldspar, and pyroxene have only 8--12 spectra each. We report balanced accuracy and cluster-bootstrap confidence intervals over rocks, but the per-class numbers for the rarest groups are uncertain, and the intervals are wide.

\paragraph{Processing differs between the two domains.} The Mars spectra come already processed by the mission pipeline (despiking, denoising, response correction) \citep{jpl2026supercamsis}, while the RRUFF spectra come from several laboratory instruments and lasers. We apply the same baseline correction and normalization to both, but some differences remain, for example in relative peak heights between 532 and 785\,nm excitation.

\paragraph{Using unlabeled Mars spectra.} The fiber-bump template is estimated from unlabeled SuperCam spectra of the development split. This is legitimate unsupervised adaptation, since no labels or test spectra are used, but it does use real flight data. The ``simulated-Mars only'' CNN, which uses no Mars data at all, is the purer lab-to-Mars result.

\paragraph{What we did not try.} Lab instruments built to mimic SuperCam \citep{manrique2024simulcam} could provide more realistic training spectra. Domain-adversarial training \citep{ganin2016domain} could use the many unlabeled SuperCam spectra in the archive. The OH band near 3690\,cm$^{-1}$ would be needed to recognize serpentine and other hydrated minerals.


\section{Conclusion}
\label{sec:conclusion}
A ``Shazam for minerals'' built from Earth laboratory spectra does not work on Mars out of the box. Allowed to choose any mineral, it decides that most Perseverance SuperCam spectra are glass. Two pieces of expert knowledge fix most of the problem: restricting predictions to minerals known from Jezero crater, and ignoring the instrument's fiber bump. Together they bring a simple library search to \balCosineWindowMars{} balanced accuracy, and to \confHighCW{} on the spectra that the SuperCam team identified with high confidence. Training a neural network on simulated Mars conditions, by contrast, made it more robust in the lab and \emph{worse} on Mars. Lab-to-Mars transfer should therefore be measured on real flight data. The benchmark we release, built from public PDS data, makes that measurement possible for anyone.
